\section{讨论与局限性}\label{sec:limitations}
\subsection{可能的机制解释}
度数加权的直觉是高度数节点拥有更多重建上下文，但它同时改变谁提供可见特征、谁承担重建目标。高度数节点更频繁地被遮蔽，可能使较多邻居失去有用的输入信号；低度数节点则因固定总预算而更少成为重建目标，可能削弱对这些节点的直接监督。均匀掩码的潜在优势是对每个节点提供近似相同的重建机会，不主动以度数重加权训练目标。

此外，SCE要求恢复原始特征方向，下游线性分类要求表示中类别信息可被简单分离，二者并非同一目标。某类节点可被邻域较容易地重建，不意味着提高其重建频率就会增强分类相关信息。重掩码、邻域相关性和标签同质性也可能共同影响结果。上述均是可能原因，本实验没有对可见邻居比例、分组重建误差或各机制分别进行消融，不能将平均下降归因于其中某一条。当前证据支持的是“掩码分布改变且分类均值下降”，而非对退化原因的因果证明。

\subsection{实验范围与协议限制}
实验仅包含Cora、五个种子及三个预设正向系数，没有跨数据集、负系数或其他权重函数的证据。以同一验证划分上的五种子均值选参仍属于有限的开发集探索，虽然测试没有反馈调参，也不提供跨数据集泛化保证。度数组的样本量和人数不均、相同测试节点在不同seed中重复使用，限制了对局部差异的统计解释。

M1与M2使用相同网络、数据和训练预算，但M2为控制策略分支的随机消耗而采用独立CPU选点。M2均匀组84.20\%与M1的83.98\%之间存在随机协议及GPU差异，不能将0.22个百分点解释为度数算法或工程适配提升。有效算法比较来自M2内部重新训练的配对控制组。

M2全部训练和前九次最终测试在L40上进行。最后一个0.25/seed4模型启动测试前，L40剩余显存不足8GiB，资源检查拒绝启动；之后补测入口的一次源码搜索路径错误也发生在任何测试前向之前。修正后在显存充足的共享A40上完成该模型唯一一次最终推断，保留失败日志和设备记录，没有重训、重选系数或重复其他九次测试。该模型的82.8\%包含不同推断设备的潜在浮点影响，应披露为限制；为避免额外测试，本文没有再次推断来估计其大小。

随机种子、cuDNN标志及RNG轨迹的一致性只能控制部分随机因素，不能保证PyG scatter等GPU归约位级确定。原论文20次初始化与本文五次、DGL与PyG历史实现、软件和硬件环境也不同。本文不强行进行显著性检验，不将一次小幅正差或平均负差解释为充分统计证据。该负结果说明预设度数策略在本次条件下没有达到改善目标，不证明所有度数策略或应用场景均无效。
